\chapter*{Prefazione}
\addcontentsline{toc}{chapter}{Prefazione}

Chernobyl, 26 aprile 1986, ore 1:23 del mattino. Nella sala di controllo del reattore numero 4, l'ingegnere Anatoly Dyatlov ordina di proseguire con un test di sicurezza nonostante tutti gli indicatori siano fuori parametro. I tecnici protestano: la potenza del reattore è scesa troppo, il sistema è instabile, bisogna fermare tutto. Dyatlov li ignora. Non è uno stupido , è un ingegnere nucleare brillante, vice capo ingegnere della centrale, con vent'anni di esperienza. Ma quella notte è convinto di tre cose: il reattore RBMK è il più sicuro del mondo, i computer sbagliano, e lui sa meglio dei giovani tecnici cosa si può fare. Ordina di estrarre troppe barre di controllo. Alle 1:23, il reattore esplode. La nube radioattiva contaminerà mezza Europa. Nei mesi successivi, durante le indagini, Dyatlov continuerà a sostenere di aver fatto tutto correttamente, che il reattore aveva un difetto nascosto, che lui non poteva saperlo. E qui sta il dettaglio più inquietante: probabilmente ci credeva davvero. Il suo cervello aveva riscritto i ricordi di quella notte per proteggere un'immagine di sé che non poteva permettersi di mettere in discussione.

Questo libro nasce da quella distorsione. Perché è così difficile ammettere i propri errori, anche di fronte a prove schiaccianti? Perché persone intelligenti e preparate ignorano segnali d'allarme evidenti? Dyatlov non era in malafede, era intrappolato in un meccanismo che le neuroscienze hanno mappato con precisione: quando la nostra identità professionale è minacciata, il cervello attiva difese cognitive potentissime. Riscrive ricordi, reinterpreta fatti, costruisce narrazioni alternative. Non è disonestà consapevole: è il modo in cui il sistema nervoso protegge il senso di sé da informazioni che potrebbero frantumarlo. Questo vale per ingegneri nucleari come per chiunque affronti decisioni ad alto rischio: medici che non vedono errori diagnostici, piloti che ignorano allarmi, manager che negano segnali di crisi. Il cervello non è progettato per cercare la verità quando la verità mette in discussione chi siamo. È progettato per mantenere coerenza interna, anche a costo di distorcere la realtà. Comprendere questo meccanismo , come funziona, quando si attiva, perché è così resistente alle evidenze , non serve solo a capire disastri storici. Serve a riconoscere quante volte, ogni giorno, anche noi riscriviamo la realtà per proteggere l'immagine che abbiamo di noi stessi.
%, , 

La neuroscienza ci ha insegnato che il cervello è una macchina straordinaria, forgiata da milioni di anni di evoluzione per risolvere problemi che i nostri antenati affrontavano nelle savane africane. Ma noi non viviamo più nelle savane. Viviamo in metropoli iperconnesse, navighiamo oceani di informazioni digitali, prendiamo decisioni che gli ominidi del Pleistocene non avrebbero neppure potuto immaginare. E qui sorge il paradosso: la nostra mente, questo prodigioso strumento di adattamento, spesso sembra inadeguata proprio nel contesto che noi stessi abbiamo creato.

Questo testo non nasce dall'ambizione di fornire un manuale d'uso completo della mente umana , impresa tanto titanica quanto illusoria. Nasce piuttosto dalla necessità di dare forma concreta a intuizioni frammentarie, di connettere puntini che apparivano sparsi in discipline diverse, di tradurre in comprensione ciò che per troppo tempo è rimasto al livello di vaga sensazione. È il tentativo di rispondere a domande che probabilmente chiunque si è posto almeno una volta: perché faccio quello che faccio? Perché gli altri agiscono in modi che mi sembrano inspiegabili? Esiste un metodo nella follia quotidiana dei nostri comportamenti?

La risposta, anticipiamolo subito, non è semplice. Ma è affascinante. E soprattutto, è fondata. Negli ultimi decenni, la ricerca neuroscientifica ha compiuto progressi che rasentano il rivoluzionario. Abbiamo imparato a guardare dentro il cervello mentre pensa, decide, ricorda, dimentica. Abbiamo scoperto che molti dei nostri comportamenti più sconcertanti non sono anomalie o debolezze caratteriali, ma conseguenze prevedibili di come il nostro sistema nervoso elabora informazioni in un mondo che cambia a una velocità vertiginosa.

Ed è qui che l'intelligenza artificiale entra in scena, non come protagonista, ma come uno specchio particolare , uno specchio alieno, potremmo dire. Nel tentativo di costruire sistemi che imitino l'intelligenza umana, abbiamo finito per creare qualcosa di radicalmente diverso da noi, e proprio questa diversità è preziosa. L'AI ci consente di guardarci da fuori, con gli occhi di chi non appartiene al nostro mondo. Riconosce i nostri pattern, replica i nostri bias cognitivi, riproduce persino i nostri errori sistematici , ma lo fa da una prospettiva esterna, come un antropologo che osserva una cultura aliena. E qui avviene qualcosa di straordinario: comportamenti che dentro di noi sembrano naturali, inevitabili, quasi invisibili, diventano chiaramente visibili quando li guardiamo riflessi in questa tecnologia estranea. La logica dei nostri errori, le tracce dei nostri pregiudizi, i pattern che guidano le nostre scelte , tutto questo emerge con nitidezza quando qualcuno (o qualcosa) che non condivide la nostra mente ci osserva e ci imita.

L'intelligenza artificiale non ci ha solo fornito nuovi strumenti tecnologici. Ci ha posto di fronte a domande antiche con una chiarezza nuova: cosa significa pensare? Come prendiamo decisioni? Perché certi comportamenti persistono nonostante la loro evidente disfunzionalità? E, forse la domanda più interessante è: quanto del nostro comportamento è davvero frutto di scelte consapevoli, e quanto invece è il risultato automatico di processi che sfuggono alla nostra percezione?

Questo libro vuole essere un ponte. Un ponte tra la ricerca scientifica e l'esperienza quotidiana, tra i laboratori di neuroscienze cognitive e le situazioni concrete che tutti viviamo ogni giorno. Vuole mostrare come certi pattern comportamentali , quelli che ci hanno sempre lasciati perplessi, quelli che abbiamo osservato ripetersi con regolarità quasi matematica senza mai comprenderne il motivo , abbiano radici profonde nella nostra architettura neurale e nel modo in cui questa interagisce con l'ambiente digitale contemporaneo.

Non troverete qui formule magiche o soluzioni rapide. Ma troverete, spero, quello che ho trovato io nel corso di questa esplorazione: una comprensione più profonda di ciò che significa essere umani nell'era dell'informazione e dell'intelligenza artificiale. Una comprensione che non elimina i problemi, ma li colloca in una cornice che li rende meno misteriosi e, paradossalmente, più gestibili.

Perché comprendere la mente non significa solo studiare neuroni e sinapsi. Significa comprendere noi stessi in tutta la nostra complessa, contraddittoria, magnificamente imperfetta umanità , un'umanità che oggi, più che mai, si riflette e si interroga attraverso gli strumenti che ha creato.

Benvenuti in questo viaggio. Un viaggio che inizia con una domanda semplice, "perché?" , e che ci porterà molto più lontano di quanto avremmo potuto immaginare.

\vspace{1cm}

\noindent \textit{L'Autore}